% Drop-in text for the PSF and Noise & Artifacts subsections.
% Figures: paper_figures/drafts/*.pdf (built by paper_figures/build_psf_noise_drafts.py).
% Numbers checked 2026-10-27 against data/euclid_psf/euclid_psf_clusters.json
% (ePSF header FWHMs) and data/euclid_stars/stars.csv on FASRC (43,401 stars).
% New citation needed: Lupton2004 (Lupton et al. 2004, PASP 116, 133) for the RGB stretch.

\subsubsection{Measuring the PSF}
To simulate Euclid observations of a synthetic sky we first need the Euclid point-spread function (PSF).
The PSF describes how Euclid blurs a point source; it is shaped by diffraction (including the spikes from the spider of the secondary mirror), optical aberrations, detector effects and pointing jitter.
Rather than adopting an instrument model, we measure the effective PSF (ePSF; \citealt{AndersonKing2000}) directly from bright point sources in the Q1 MER mosaics with \texttt{photutils}.
Importantly, the ePSF also captures dithering and resampling, which matter most for NISP, whose $0.3''$ native pixels are resampled onto the $0.1''$ MER mosaic grid (Fig.~\ref{fig:psfs}, bottom row).
We select $\sim$43{,}000 unmasked point-like sources with $16 < \mathrm{VIS} < 19$ from the MER catalogue, extract $51''$ cutouts in all four bands, and reject sources with saturated cores, keeping 17{,}905 stars that are usable in every band (Fig.~\ref{fig:psf_measurement}).
To capture spatial variation we cluster them by sky position into 44 regions of 270--590 stars each (median 397) and build one ePSF per region and band, oversampled by a factor of two ($0.05''$ pixels; Fig.~\ref{fig:psf_variation}).
Finally, a radial cosine taper takes each stamp smoothly to zero between $21''$ and $25''$ from its centre, so no square stamp edge is imprinted around bright stars (Fig.~\ref{fig:psf_profiles}c).
Over the 44 regions the median FWHMs are $0.16''$ (VIS), $0.36''$ ($Y_E$), $0.36''$ ($J_E$) and $0.39''$ ($H_E$).%
% Medians exclude 7 failed fits (6 in Y_E, 1 in H_E) at 0.08--0.10'', below the NISP
% diffraction limit. Including them: 0.16 / 0.35 / 0.36 / 0.39''.

\begin{figure*}
    \centering
    \includegraphics[width=\linewidth]{images/psf_gallery.pdf}
    \caption{Empirical ePSFs of one region (EDF-S, 496 stars) in the four Euclid bands.
    \emph{Top:} the central $25''$ on a logarithmic stretch, normalised to the peak, showing the six diffraction spikes and, in $H_E$, the first Airy ring.
    \emph{Bottom:} the central $2''$ on a linear stretch at the $0.05''$ ePSF sampling, with the native detector pixel grid overlaid ($0.1''$ for VIS, $0.3''$ for NISP) and a dashed circle of the median FWHM over all 44 regions.
    The NISP core spans only a few native pixels, so its ePSF also encodes the dithering and resampling of the MER mosaics.}
    \label{fig:psfs}
\end{figure*}

\begin{figure*}
    \centering
    \includegraphics[width=\linewidth]{images/psf_measurement.pdf}
    \caption{From individual stars to the ePSF.
    Each row shows five Q1 stars ($16<\mathrm{VIS}<19$) in one band, normalised to their peak on a logarithmic stretch ($8''$ stamps).
    Stars like these are aligned, oversampled by a factor of two and stacked into one ePSF per region; the right column shows the ePSF of the region in Fig.~\ref{fig:psfs}, built from 496 stars.
    The stacked kernel recovers the wings that are lost in the noise of any single star.}
    \label{fig:psf_measurement}
\end{figure*}

\begin{figure}
    \centering
    \includegraphics[width=\linewidth]{images/psf_rgb.pdf}
    \caption{The chromatic Euclid PSF: the $H_E$, $Y_E{+}J_E$ and VIS ePSFs of Fig.~\ref{fig:psfs}, each normalised to unit flux, shown as red, green and blue with a colour-preserving arcsinh stretch \citep{Lupton2004}.
    \emph{Left:} $25''$; \emph{right:} the central $2''$ (box).
    The sharp VIS core sits inside the broader NISP halo.
    A $0.05''$ Gaussian smoothing is applied for display.}
    \label{fig:psf_rgb}
\end{figure}

\begin{figure*}
    \centering
    \includegraphics[width=\linewidth]{images/psf_variation.pdf}
    \caption{Spatial variation of the PSF.
    \emph{Top left:} the 44 star regions in the three Q1 fields; markers sit at the region centroids (area proportional to the number of stars), cells show the nearest-centroid partition, and colour gives the VIS FWHM.
    \emph{Top right:} FWHM of every region in each band; bars mark the medians.
    Open symbols are fits below the NISP diffraction limit and are excluded from the medians quoted in the text.
    \emph{Bottom:} the VIS ePSF of every region ($12''$, logarithmic stretch), grouped by field and framed by its FWHM colour.
    Many regions show more than six spikes, likely because they combine stars observed at different roll angles.}
    \label{fig:psf_variation}
\end{figure*}

\begin{figure*}
    \centering
    \includegraphics[width=\linewidth]{images/psf_profiles.pdf}
    \caption{Radial structure of the ePSFs.
    (a) Azimuthally averaged profile of the mean ePSF in each band, normalised to the peak; shading spans the 16th--84th percentiles over the 44 regions, and dashed lines show the profiles before the edge taper (grey band).
    Beyond $\sim$10$''$ the profiles flatten at $\sim$$2\times10^{-6}$ (VIS) and $\sim$$2\times10^{-5}$ (NISP) of the peak.
    (b) Encircled energy of the tapered mean ePSFs.
    (c) The radial cosine taper on the VIS ePSF at a deep stretch: the raw stamp (left half) carries a square edge, the tapered stamp (right half) does not. Circles mark where the taper begins and ends.}
    \label{fig:psf_profiles}
\end{figure*}

\subsubsection{Noise \& Artifacts}
% (paragraph still to be written)

\begin{figure*}
    \centering
    \includegraphics[width=\linewidth]{images/noise_artifacts.pdf}
    \caption{The noise and artifact model.
    \emph{Top:} an illustrative $25.6''$ scene (TNG50 galaxies and stars) taken from (a) the $0.05''$ truth through (b) convolution with the ePSFs and binning onto the $0.1''$ MER grid, (c) sky and source photon noise with the correlated texture of a dithered stack, (d) a pointing seam and the wing of a bright star outside the field, to (e) surviving detector artifacts and saturation blackout.
    The composites show $H_E$, $Y_E{+}J_E$ and VIS as red, green and blue, each in units of its sky $\sigma$.
    \emph{Bottom:} (f) unit-variance noise: white, and the stacked texture in VIS and $H_E$, where each $0.3''$ detector pixel is shared among $3\times3$ mosaic pixels;
    (g) local VIS noise $\sigma$ of panel (d) relative to the sky, raised over bright sources by their photon noise and by a factor 1.38 across the seam (the off-field wing also shows);
    (h) distribution of the sky noise measured in the Q1 MER noise maps at 294 tile positions, relative to each band's median (log scatter on the right);
    (i) a $6''$ VIS zoom with cosmic-ray hits, hot and dead pixels and a masked trail marked;
    (j) a simulated saturation blackout of a bright star (top) and a galaxy nucleus in Q1 that is blanked in $H_E$ but intact in VIS (bottom).
    For visibility the demo scene forces the seam, the wing and the blackout, and boosts the rates of cosmic rays ($\times6$), hot pixels ($\times60$) and trails ($\times25$) over those used to generate the training data.}
    \label{fig:noiseandartifacts}
\end{figure*}
